\section{Betti bounds for stable ideals}
\label{sec:tor}

The preceding section shows that degreewise lex replacement of a
monomial ideal stable in the box is again an ideal with the same
Hilbert function.  We prove that this replacement also bounds every
graded Betti number.  The induction on the number of variables uses
the first and last box slices.  After reduction modulo the last variable,
these determine two endpoint modules; the middle modules have zero
Koszul differentials.

For modules with the same Hilbert function, the Koszul terms have equal
dimensions in every homological and internal degree.  An upper bound on
Betti numbers is therefore equivalent to a lower bound on the sum of the
two adjacent differential ranks.  We first isolate this rank sum, which
does not increase under taking either submodules or quotients.
Let $m\ge0$ and $R=\C[x_1,\ldots,x_m]$.  For a graded $R$-module $B$ with
finite-dimensional graded pieces, write $K^R_\bullet(B)$ for its
Koszul complex and define its \emph{Koszul rank deficit} by
\begin{equation}
\label{tor:deficit-definition}
 D^R_{p,j}(B)=\dim_\C K^R_p(B)_j-\beta^R_{p,j}(B),
 \qquad p\ge0,\quad j\in\Z.
\end{equation}
Terms of the Koszul complex outside its homological range are zero.

\begin{lemma}
\label{tor:deficit}
For every $p\ge0$ and $j\in\Z$, the quantity $D^R_{p,j}(B)$
is the sum of the ranks of the two Koszul differentials adjacent to
$K^R_p(B)_j$.  If $U\subseteq B$ is a graded submodule and $V$ is
a graded quotient of $B$, then
\[
 D^R_{p,j}(U)\le D^R_{p,j}(B),\qquad
 D^R_{p,j}(V)\le D^R_{p,j}(B).
\]
The quantity is additive on finite direct sums, satisfies
\[
 D^R_{p,j}(B(-s))=D^R_{p,j-s}(B),
\]
and vanishes when $R_{>0}B=0$.
\end{lemma}

\begin{proof}
Write $d_p(B)_j:K^R_p(B)_j\to K^R_{p-1}(B)_j$ for the Koszul
differential.  Since its homology computes $\Tor^R(B,\C)$,
\[
 \beta^R_{p,j}(B)
 =\dim_\C K^R_p(B)_j
       -\operatorname{rank}d_p(B)_j
       -\operatorname{rank}d_{p+1}(B)_j.
\]
This proves the rank formula.  The Koszul complex of $U$ embeds
termwise into that of $B$, and the image of each of its differentials
embeds into the corresponding differential image for $B$.
For the quotient $V$, the termwise surjections give
\[
 \operatorname{im}d_p(V)_j
   =\text{the image of }\operatorname{im}d_p(B)_j
       \text{ in }K^R_{p-1}(V)_j.
\]
Thus each of the two ranks can only decrease under either operation.
Direct sums and grading shifts commute with the Koszul complex.
Finally, if all the variables annihilate $B$, every Koszul differential
is zero.
\end{proof}

\begin{theorem}
\label{tor:stable}
Let $r\ge1$, let $2\le a_1\le\cdots\le a_r$, and let
$A=\C[x_1,\ldots,x_r]$.  Suppose that $H\subseteq A$ is a
monomial ideal containing $(x_1^{a_1},\ldots,x_r^{a_r})$ and is
stable in the box.  Let $G=\Lambda H$ be its degreewise lex
replacement inside the box, together with the prescribed pure powers.
Then, for every $p\ge0$ and $j\in\Z$,
\[
 \beta^A_{p,j}(H)\le\beta^A_{p,j}(G),
 \qquad
 \beta^A_{p,j}(A/H)\le\beta^A_{p,j}(A/G).
\]
\end{theorem}

\begin{proof}
By Proposition~\ref{box:ideal}, $G$ is an ideal stable in the box
and has the same Hilbert function as $H$.  If $H=A$, then $G=A$
and both comparisons are equalities.  For proper $H$, the ideal
comparison implies the quotient comparison in positive homological
degrees by Lemma~\ref{filt:minimal}:
\begin{equation}
\label{tor:ideal-quotient}
 \beta^A_{p+1,j}(A/H)=\beta^A_{p,j}(H),
 \qquad
 \beta^A_{p+1,j}(A/G)=\beta^A_{p,j}(G).
\end{equation}
The two quotient Betti numbers in homological degree zero are both
$1$ when $j=0$ and $0$ otherwise.  We prove the two comparisons
simultaneously by induction on $r$, establishing the ideal comparison
and then using~\eqref{tor:ideal-quotient}.  When $r=1$, a monomial
ideal containing $x_1^{a_1}$ is a power ideal or the full ring, and
$H=G$.

Assume $r>1$ and $H$ proper, and put
\[
 R=\C[x_1,\ldots,x_{r-1}],\qquad z=x_r,\qquad a=a_r.
\]
For $s\ge0$, define coefficient ideals
\[
 H^{(s)}=\{u\in R:uz^s\in H\},\qquad
 G^{(s)}=\{u\in R:uz^s\in G\}.
\]
These are ascending monomial ideals containing the earlier prescribed
powers, and they equal $R$ for $s\ge a$.  For $s<a$, their
stability in the earlier-variable box follows by applying stability
with the last exponent fixed at $s$.  For $1\le s<a$ we also have
\begin{equation}
\label{tor:slice-annihilation}
 R_{>0}H^{(s)}\subseteq H^{(s-1)},\qquad
 R_{>0}G^{(s)}\subseteq G^{(s-1)}.
\end{equation}
To see the first inclusion, it is enough to multiply a monomial of
$H^{(s)}$ by an earlier variable $x_i$.  If its earlier-variable
exponent, or that exponent after multiplication by $x_i$, is outside
the box, the earlier powers already imply the required membership.
Otherwise stability moves one factor of $z$ to $x_i$.  The proof
for $G$ is identical.

We next pass from ideals over $A$ to modules over $R$.  Multiplication
by $z$ is injective on $H$, since $H$ is an ideal in the domain $A$.
Set
\[
 B_H=H/zH,\qquad B_G=G/zG.
\]
Reducing a minimal $A$-free resolution of $H$ modulo $z$ gives a
minimal $R$-free resolution of $B_H$.  Here is a direct exactness
check.  Lift a cycle in positive homological degree modulo $z$ to
an element $y$ of the free resolution, and write $dy=zv$.
Applying the next differential shows that $v$ is a cycle, using
injectivity of $z$ on the free modules; in degree zero, apply the
augmentation and use injectivity of $z$ on $H$ instead.
Exactness gives $v=du$, so $y-zu$ is a cycle and hence a boundary.
Its reduction modulo $z$ is the original cycle.  Right exactness
identifies the cokernel with $B_H$, and all differential entries
remain of positive degree or become zero, so minimality is preserved.
The same argument applies to $G$.  Therefore
\begin{equation}
\label{tor:reduction}
 \beta^A_{p,j}(H)=\beta^R_{p,j}(B_H),\qquad
 \beta^A_{p,j}(G)=\beta^R_{p,j}(B_G).
\end{equation}
The injective multiplication maps also give
\[
 \dim_\C(B_H)_d=\dim_\C H_d-\dim_\C H_{d-1},\qquad
 \dim_\C(B_G)_d=\dim_\C G_d-\dim_\C G_{d-1}.
\]
Consequently $B_H$ and $B_G$ have equal Hilbert functions.

The monomial decomposition by powers of $z$ describes these modules
more precisely; this is the decomposition appearing in
\cite[Lemma~3.7]{CS}.  The coefficient of $z^0$ in $H/zH$ is $H^{(0)}$,
and its coefficient of $z^s$ for $s\ge1$ is
$H^{(s)}/H^{(s-1)}$.  Since $H^{(s)}=R$ for $s\ge a$, this gives
an isomorphism of graded $R$-modules
\begin{equation}
\label{tor:slice-decomposition}
 B_H\cong H^{(0)}
       \oplus\bigoplus_{s=1}^{a-1}
             (H^{(s)}/H^{(s-1)})(-s)
       \oplus(R/H^{(a-1)})(-a).
\end{equation}
There are no terms for $s>a$.  The shifts record that multiplication
by $z^s$ raises ordinary degree by $s$.  By
\eqref{tor:slice-annihilation}, all middle summands are annihilated
by the variables of $R$, as illustrated in
Figure~\ref{fig:endpoint-slices}. The decomposition is $R$-linear
because multiplying by an earlier variable preserves the exponent of
$z$. The middle summands still contribute to $\Tor$, but their
Koszul differentials have rank zero. Thus they contribute nothing to
the deficits. Lemma~\ref{tor:deficit} yields
\begin{equation}
\label{tor:endpoint-deficits}
 D^R_{p,j}(B_H)
   =D^R_{p,j}(H^{(0)})
      +D^R_{p,j-a}(R/H^{(a-1)}),
\end{equation}
and the identical formula holds for $G$.

\begin{figure}[htbp]
\centering
\begin{tikzpicture}[>=Stealth,font=\small]
\node (u) at (-3.3,0) {$uz^s\in H$};
\node (xu) at (3.3,0) {$x_iuz^s\in zH$};
\node (transfer) at (3.3,-1.5) {$x_iuz^{s-1}\in H$};
\draw[->] (u) -- node[above] {$\times x_i$} (xu);
\draw[->] (transfer) -- node[right] {$\times z$} (xu);
\draw[->,dashed] (u) --
  node[pos=.47,below,sloped] {transfer $z$ to $x_i$} (transfer);
\node[align=center] at (-3.2,-1.55)
  {$x_i[uz^s]=0$ in $H/zH$\\[3pt]
   $1\le s<a$};
\end{tikzpicture}
\caption{A middle slice is annihilated by every earlier variable.
The dashed arrow is supplied by box stability, or by an earlier
prescribed power if the monomial leaves the earlier-variable box.
Consequently all Koszul differentials on a middle slice vanish.
At $s=0$ there is no factor of $z$ to transfer; at $s=a$ the
last exponent is outside the box. These are the two endpoints.}
\label{fig:endpoint-slices}
\end{figure}


It remains to compare the two endpoint contributions.  Write
$\Lambda'$ for lex replacement with the earlier prescribed powers
in $R$.  The induction hypothesis applies to both $H^{(0)}$ and
$H^{(a-1)}$.  It gives an ideal comparison for the first and a
quotient comparison for the second.  Their respective lex
replacements have the same Hilbert functions, and hence the same
Koszul term dimensions.  Subtracting the Betti inequalities from
these dimensions gives
\begin{align}
 D^R_{p,j}(H^{(0)})
       &\ge D^R_{p,j}(\Lambda'H^{(0)}),
       \label{tor:induction-zero}\\
 D^R_{p,j-a}(R/H^{(a-1)})
       &\ge D^R_{p,j-a}(R/\Lambda'H^{(a-1)}).
       \label{tor:induction-top}
\end{align}
This also covers any coefficient ideal equal to $R$: its lex
replacement is $R$, the corresponding ideal comparison is equality,
and its quotient is zero.

The slice inequalities of Lemma~\ref{box:inequalities} imply
\begin{equation}
\label{tor:endpoint-inclusions}
 G^{(0)}\subseteq\Lambda'H^{(0)},\qquad
 \Lambda'H^{(a-1)}\subseteq G^{(a-1)}.
\end{equation}
Indeed, fix an earlier-variable degree $e\ge0$.  For the first
inclusion, apply the zero-slice inequality to the included exponent
sets of $H$ and $G$ in total degree $e$.  For the second, apply
the top-slice inequality in total degree $e+a-1$.  Each slice of
$G$ is already a lex segment in the earlier-variable box, so its
cardinality comparison with the corresponding slice of $H$ becomes
the stated containment in the appropriate lex replacement.
All monomials outside the earlier-variable box belong to both ideals.
This proves the inclusions in every degree.

Apply Lemma~\ref{tor:deficit} to the first inclusion in
\eqref{tor:endpoint-inclusions} and to the quotient map furnished
by the second:
\[
 R/\Lambda'H^{(a-1)}\longrightarrow R/G^{(a-1)}.
\]
Together with~\eqref{tor:induction-zero} and
\eqref{tor:induction-top}, this gives
\begin{align*}
 D^R_{p,j}(H^{(0)})&\ge D^R_{p,j}(G^{(0)}),\\
 D^R_{p,j-a}(R/H^{(a-1)})
       &\ge D^R_{p,j-a}(R/G^{(a-1)}).
\end{align*}
Adding and using~\eqref{tor:endpoint-deficits}, we obtain
\[
 D^R_{p,j}(B_H)\ge D^R_{p,j}(B_G).
\]
The equal Hilbert functions of $B_H$ and $B_G$ imply
\[
 \dim_\C K^R_p(B_H)_j
    =\binom{r-1}{p}\dim_\C(B_H)_{j-p}
    =\dim_\C K^R_p(B_G)_j,
\]
with the binomial coefficient interpreted as zero outside its range.
Subtracting the deficit inequality from these equal dimensions gives
$\beta^R_{p,j}(B_H)\le\beta^R_{p,j}(B_G)$.
Equation~\eqref{tor:reduction} proves the ideal comparison over $A$,
and~\eqref{tor:ideal-quotient} completes the quotient comparison.
Every step applies to every $p\ge0$ and $j\in\Z$; in particular,
negative internal degrees and Koszul terms outside their homological
range introduce no additional cases.
\end{proof}
